\chapter{Microstructure du change au comptant}\label{ch:m2:fx-spot-microstructure}

Le 15 décembre 2010, un client de la plateforme de change électronique d'une grande banque lui
écrit~: «~nous avons remarqué qu'il y avait plus de 300 ordres rejetés par vous
aujourd'hui et la raison est `NACK'. Pourriez-vous y jeter un œil, s'il vous plaît.» Deux jours
plus tard, sans réponse, il réécrit~: «~nous continuions de recevoir les meilleurs cours affichés
par vous et à toucher votre cours, mais nous avons été rejetés 9 fois sur
10~». Le New York State Department of Financial Services cite les deux messages
cinq ans plus tard, dans l'ordonnance de consentement par laquelle Barclays a versé
150~millions de dollars, sans trouver de preuve que la banque ait jamais répondu. Les rejets
provenaient d'une pratique appelée \omterm{def:m2:fx-spot-microstructure:lastlook}{last look}~: la banque retenait chaque ordre pendant quelques dizaines ou
centaines de millisecondes, et le refusait si le cours avait évolué à son détriment
entre-temps. Ce chapitre porte sur la mécanique du change électronique
au comptant~: comment les cours sont diffusés en streaming aux clients, pourquoi le cours obtenu
dépend de qui il est, comment fonctionne le \omterm{def:m2:fx-spot-microstructure:lastlook}{last look} et ce qu'il coûte
(\cref{fig:m2:fx-spot-microstructure:timeline}), comment les dealers décalent leurs cotations et
internalisent les flux, et le code de conduite que le marché s'est lui-même donné.

\section{Cotations en streaming et liquidité par paliers}

\begin{definition}[Cotation en streaming, palier de liquidité]\label{def:m2:fx-spot-microstructure:stream}
Une \emph{cotation en streaming}\index{cotation en streaming} est un cours acheteur et un cours vendeur, pour des tailles données,
qu'un fournisseur de liquidité publie en continu à un client ou à une
plateforme et met à jour plusieurs fois par seconde~; le client négocie en envoyant une
demande au prix affiché. Un \emph{palier de liquidité}\index{palier de liquidité} est
l'ensemble des fourchettes et des tailles qu'un fournisseur affiche à un groupe de clients, fixées selon
sa perception de leurs flux.
\end{definition}

Le change au comptant électronique se négocie surtout en streaming. Le moteur de tarification d'un dealer
reçoit les \omterm{def:m2:the-fx-market:venue}{plateformes principales}, d'autres plateformes et son propre flux, et se forge une idée du
mid, et publie à chaque client un cours autour de ce niveau~: serré pour les clients dont
les transactions précèdent rarement les mouvements de prix, plus large pour les autres, et parfois
encore plus large lorsque les marchés vont vite. Les sociétés de trading agissant en principal diffusent elles aussi
directement aux banques et sur les plateformes (au \cref{ch:m2:the-fx-market}). Un client qui accède à
de nombreux flux voit le meilleur cours acheteur et le meilleur cours vendeur agrégés, et envoie sa
demande à celui qui affiche le meilleur prix.

Le problème du dealer est qu'une \omterm{def:m2:fx-spot-microstructure:stream}{cotation en streaming} est une option gratuite. Elle affiche un prix pendant
quelques millisecondes, et quiconque sait que le marché a bougé avant que la cotation
ait été mise à jour peut négocier du côté périmé. La fourchette est la protection du
dealer contre le bruit des flux non informés~; les paliers, le \omterm{def:m2:fx-spot-microstructure:lastlook}{last look} et les décalages de cotations sont
ses protections contre le reste.

\section{Sélection adverse}

Les pires clients d'un teneur de marché sont ceux qui négocient juste avant que le cours ne bouge
dans leur sens, c'est-à-dire la sélection adverse du livre~1, chapitre~1. Sur le change au comptant, la
forme la plus pure est l'arbitrage de latence~: une société rapide voit le cours bouger sur une plateforme
et frappe les cotations périmées des fournisseurs plus lents ailleurs avant qu'ils ne les mettent à jour. Un
fournisseur mesure ses clients au moyen de \omterm{def:m2:fx-spot-microstructure:markout}{mark-outs}.

\begin{definition}[Mark-out]\label{def:m2:fx-spot-microstructure:markout}
Le \emph{mark-out}\index{mark-out} d'une transaction à l'horizon $H$ est le profit
réalisé par le fournisseur en valorisant sa position au mid $H$ après la transaction~: la
demi-fourchette perçue, moins le mouvement du mid en faveur du client sur $H$.
Moyenné sur les transactions d'un client à plusieurs horizons, de la milliseconde à la
minute, il mesure à quel point le flux du client est informé.
\end{definition}

Un client dont les \omterm{def:m2:fx-spot-microstructure:markout}{mark-outs} sont proches de la demi-fourchette à tous les horizons est
non informé~; son flux a de la valeur, et il obtient le meilleur palier. Un client dont les
\omterm{def:m2:fx-spot-microstructure:markout}{mark-outs} deviennent négatifs en moins d'une seconde négocie sur une information que le
fournisseur ne possédait pas~; on lui affiche des cours plus larges, des confirmations plus lentes, ou rien.

\section{Last look et temps de rétention}

\begin{definition}[Last look, temps de rétention, taux de rejet]\label{def:m2:fx-spot-microstructure:lastlook}
Le \emph{last look}\index{last look} est une pratique par laquelle un fournisseur de liquidité
recevant une demande de transaction à son cours affiché dispose d'une dernière possibilité, une courte
fenêtre après l'arrivée de la demande, pour l'accepter ou la rejeter. Le \emph{temps de
rétention}\index{temps de rétention} est la durée de cette fenêtre~; le \emph{taux de
rejet}\index{taux de rejet} est la part des demandes rejetées. Le contrôle de prix peut
être \emph{asymétrique}, ne rejetant que les demandes dont le cours a évolué dans le
sens favorable au client au-delà d'un seuil, ou \emph{symétrique}, rejetant les mouvements au-delà
de ce seuil, dans un sens comme dans l'autre.
\end{definition}

L'argument du fournisseur en faveur du \omterm{def:m2:fx-spot-microstructure:lastlook}{last look} est qu'il peut coter plus serré s'il peut
refuser les transactions qui arrivent sur des cours périmés. L'objection du client est qu'une
demande rejetée le laisse avec le risque de marché dont il voulait se délester, à un moins bon
cours, et qu'un contrôle asymétrique permet au fournisseur de garder les transactions dont
le cours a évolué dans son sens, et de refuser les autres.

\begin{omfigure}
\centering
\begin{tikzpicture}[font=\small,x=1.3cm]
\draw[thick,-{Stealth}] (0,0) -- (10.4,0) node[right,font=\footnotesize] {temps};
\foreach \x/\l in {0/cotation\\envoyée,2.2/arrivée de\\la demande,6.6/décision} {\draw (\x,0.12) -- (\x,-0.12); \node[below,font=\footnotesize,align=center] at (\x,-0.15) {\l};}
\fill[omThm!15] (2.2,0.2) rectangle (6.6,1.0);
\node[omThm,font=\footnotesize,align=center] at (4.4,0.6) {temps de rétention~:\\le fournisseur conserve l'option};
\draw[omDef,thick,-{Stealth}] (6.6,1.25) -- (7.6,1.85) node[right,font=\footnotesize,align=left] {accepter~:\\transaction au cours affiché};
\draw[omProp,thick,-{Stealth}] (6.6,1.25) -- (7.6,0.65) node[right,font=\footnotesize,align=left] {rejeter~:\\le client conserve son risque};
\draw[gray,decorate,decoration={brace,amplitude=4pt,mirror}] (0,-0.9) -- (2.2,-0.9);
\node[gray,font=\footnotesize] at (1.1,-1.3) {latence};
\end{tikzpicture}

\omcaption{Le parcours d'une demande sous \omterm{def:m2:fx-spot-microstructure:lastlook}{last look}. La cotation parvient au client,
la demande du client parvient au fournisseur après une latence, et le fournisseur
décide à la fin du \omterm{def:m2:fx-spot-microstructure:lastlook}{temps de rétention} s'il conclut la transaction au cours qu'il a affiché.
Pendant ce \omterm{def:m2:fx-spot-microstructure:lastlook}{temps de rétention}, c'est le fournisseur, et non le client, qui détient le choix. Schéma de principe.}\label{fig:m2:fx-spot-microstructure:timeline}
\end{omfigure}

\begin{proposition}[Ce que prend une fenêtre asymétrique]\label{prop:m2:fx-spot-microstructure:transfer}
Supposons que le cours évolue en faveur du client pendant le \omterm{def:m2:fx-spot-microstructure:lastlook}{temps de rétention} $h$, d'un
montant $X \sim N(0, \sigma^2 h)$, indépendant de la demande, comme c'est le cas pour un
client non informé. Un contrôle asymétrique de seuil $\theta$ rejette lorsque
$X > \theta$, et le client abandonne, par demande,
\[
\mathbb{E}\bigl[X\,\mathbf 1\{X > \theta\}\bigr] \;=\; \sigma\sqrt{h}\;\varphi\!\left(\frac{\theta}{\sigma\sqrt h}\right).
\]
Un contrôle symétrique rejette $|X| > \theta$ et ne coûte rien en moyenne
(la \cref{fig:m2:fx-spot-microstructure:tails}).
\end{proposition}

\begin{proof}
Pour $X = sZ$ avec $s = \sigma\sqrt h$, $\mathbb{E}[sZ\mathbf 1\{Z > \theta/s\}] =
s\varphi(\theta/s)$. Sous le contrôle symétrique, la région rejetée est symétrique
par rapport à zéro, de sorte que les mouvements rejetés ont une moyenne nulle.
\end{proof}

\begin{omfigure}
\begin{tikzpicture}
\begin{axis}[name=left,width=6.4cm,height=4.4cm,
  xlabel={mouvement en faveur du client (pb)},yticklabels={},ytick=\empty,
  tick label style={font=\small},label style={font=\small},
  xmin=-0.35,xmax=0.35,ymin=0,ymax=4.4,domain=-0.35:0.35,samples=120,
  title={asymétrique},title style={font=\small}]
\addplot[fill=omProp!35,draw=none,domain=0.05:0.35] {exp(-0.5*(x/0.1)^2)/(0.1*sqrt(2*pi))} \closedcycle;
\addplot[omThm,very thick] {exp(-0.5*(x/0.1)^2)/(0.1*sqrt(2*pi))};
\draw[gray,dashed] (axis cs:0.05,0) -- (axis cs:0.05,4.2);
\node[omProp,font=\scriptsize,align=center] at (axis cs:0.22,2.6) {rejetées\\30.9~\%};
\end{axis}
\begin{axis}[at={(left.east)},anchor=west,xshift=1.0cm,width=6.4cm,height=4.4cm,
  xlabel={mouvement en faveur du client (pb)},yticklabels={},ytick=\empty,
  tick label style={font=\small},label style={font=\small},
  xmin=-0.35,xmax=0.35,ymin=0,ymax=4.4,domain=-0.35:0.35,samples=120,
  title={symétrique},title style={font=\small}]
\addplot[fill=omProp!35,draw=none,domain=0.05:0.35] {exp(-0.5*(x/0.1)^2)/(0.1*sqrt(2*pi))} \closedcycle;
\addplot[fill=omProp!35,draw=none,domain=-0.35:-0.05] {exp(-0.5*(x/0.1)^2)/(0.1*sqrt(2*pi))} \closedcycle;
\addplot[omThm,very thick] {exp(-0.5*(x/0.1)^2)/(0.1*sqrt(2*pi))};
\draw[gray,dashed] (axis cs:0.05,0) -- (axis cs:0.05,4.2);
\draw[gray,dashed] (axis cs:-0.05,0) -- (axis cs:-0.05,4.2);
\node[omProp,font=\scriptsize,align=center] at (axis cs:0.22,2.6) {rejetées\\61.7~\%};
\end{axis}
\end{tikzpicture}

\omcaption{Le mouvement du cours pendant un \omterm{def:m2:fx-spot-microstructure:lastlook}{temps de rétention} de 100 millisecondes, pour une
demande non informée (écart type de 0.1 point de base), et les demandes
rejetées avec un seuil de 0.05 point de base (en grisé). Le contrôle asymétrique ne rejette
que la queue de droite, dont le client perd la moyenne~; le contrôle symétrique rejette les deux
queues, qui se compensent. Illustratif.}\label{fig:m2:fx-spot-microstructure:tails}
\end{omfigure}

Le contrôle symétrique coûte tout de même quelque chose au client, les exécutions qu'il perd,
mais il ne sélectionne pas contre lui. C'est la distinction qu'ont retenue les régulateurs. Dans
l'affaire de l'ordre new-yorkais, Barclays avait comparé le cours au début et à la fin d'un
\omterm{def:m2:fx-spot-microstructure:lastlook}{temps de rétention} de plusieurs dizaines à plusieurs centaines de millisecondes, et rejeté les ordres
devenus défavorables pour elle~; elle avait rendu le contrôle
symétrique en septembre et octobre 2014, et une plateforme, avec 7~\% de son
volume, est restée asymétrique jusqu'en août 2015.

\begin{omfigure}
\begin{tikzpicture}
\begin{axis}[name=left,width=6.4cm,height=5cm,xmode=log,log basis x=10,
  xlabel={temps de rétention (ms)},ylabel={taux de rejet (\%)},
  tick label style={font=\small},label style={font=\small},
  xmin=8,xmax=600,ymin=0,ymax=105,grid=major,grid style={gray!25},
  xtick={10,25,50,100,200,500},xticklabels={10,25,50,100,200,500},
  legend columns=2,legend style={font=\footnotesize,at={(1.1,-0.32)},anchor=north,draw=none,fill=none,/tikz/every even column/.append style={column sep=8pt}},
  table/col sep=comma]
\addplot[omProp,very thick,mark=*,mark size=1.2pt] table[x=hold,y=rej_inf_asym]{figdata/markets-2/15-fx-spot-microstructure/hold.csv};
\addplot[omProp,thick,dashed,mark=o,mark size=1.4pt] table[x=hold,y=rej_inf_sym]{figdata/markets-2/15-fx-spot-microstructure/hold.csv};
\addplot[omDef,very thick,mark=*,mark size=1.2pt] table[x=hold,y=rej_uninf_asym]{figdata/markets-2/15-fx-spot-microstructure/hold.csv};
\addplot[omDef,thick,dashed,mark=o,mark size=1.4pt] table[x=hold,y=rej_uninf_sym]{figdata/markets-2/15-fx-spot-microstructure/hold.csv};
\legend{informé asymétrique,informé symétrique,non informé asymétrique,non informé symétrique}
\end{axis}
\begin{axis}[at={(left.east)},anchor=west,xshift=1.4cm,width=6.4cm,height=5cm,xmode=log,log basis x=10,
  xlabel={temps de rétention (ms)},ylabel={mark-out (pb)},
  tick label style={font=\small},label style={font=\small},
  xmin=8,xmax=600,ymin=0.35,ymax=0.6,grid=major,grid style={gray!25},
  xtick={10,25,50,100,200,500},xticklabels={10,25,50,100,200,500},
  table/col sep=comma]
\addplot[omThm,very thick,mark=*,mark size=1.2pt] table[x=hold,y=mo_asym]{figdata/markets-2/15-fx-spot-microstructure/hold.csv};
\addplot[omThm,thick,dashed,mark=o,mark size=1.4pt] table[x=hold,y=mo_sym]{figdata/markets-2/15-fx-spot-microstructure/hold.csv};
\draw[gray,dashed] (axis cs:8,0.3785) -- (axis cs:600,0.3785);
\node[font=\scriptsize,gray,anchor=north east] at (axis cs:590,0.3785) {pas de last look};
\end{axis}
\end{tikzpicture}

\omcaption{Un flux EURUSD simulé avec 10~\% de demandes informées. À gauche~: \omterm{def:m2:fx-spot-microstructure:lastlook}{taux de rejet}
selon le \omterm{def:m2:fx-spot-microstructure:lastlook}{temps de rétention} pour des clients informés et non informés, sous un contrôle asymétrique et
un contrôle symétrique au même seuil de 0.05 point de base. À droite~: le
\omterm{def:m2:fx-spot-microstructure:markout}{mark-out} à une seconde du fournisseur par demande exécutée (trait plein asymétrique, tirets
symétrique). Les deux contrôles arrêtent la plupart des demandes informées~; seul le contrôle
asymétrique transforme des \omterm{def:m2:fx-spot-microstructure:lastlook}{temps de rétention} plus longs en profit supplémentaire, prélevé sur les clients non informés.
Paramètres illustratifs~; données~: le tutoriel du chapitre.}\label{fig:m2:fx-spot-microstructure:hold}
\end{omfigure}

La \cref{fig:m2:fx-spot-microstructure:hold} montre les deux effets sur un flux simulé. Les demandes
informées, qui arrivent après un mouvement de 0.2 point de base que la cotation n'a pas
encore rattrapé, sont presque toutes rejetées par l'un ou l'autre contrôle pour les \omterm{def:m2:fx-spot-microstructure:lastlook}{temps
de rétention} courts. Les demandes non informées sont davantage rejetées par le contrôle
symétrique, qui refuse les mouvements dans les deux sens, mais cela ne leur coûte rien en
moyenne~; le contrôle asymétrique en rejette moins, et son \omterm{def:m2:fx-spot-microstructure:markout}{mark-out} croît avec le \omterm{def:m2:fx-spot-microstructure:lastlook}{temps de
rétention}, car chaque milliseconde supplémentaire donne au cours plus de temps pour bouger
en faveur du client et être refusé.

\begin{example}[Le coût d'une fenêtre]\label{ex:m2:fx-spot-microstructure:window}
Avec la volatilité du flux, 0.1 point de base sur 100 millisecondes, un contrôle
asymétrique de seuil 0.05 point de base et un \omterm{def:m2:fx-spot-microstructure:lastlook}{temps de rétention} de 100 millisecondes
prend $0.1\,\varphi(0.5) = 0.0352$ point de base par demande aux clients non informés~; sans
seuil, 0.0399. Sur 2~milliards de dollars de demandes par jour, cela fait
environ 7\,041~dollars par jour, soit 1.76~million de dollars sur 250 jours de négociation. Le
contrôle symétrique au même seuil rejette 62~\% des demandes et ne coûte rien en
moyenne.
\end{example}

\section{Décalage des cotations (skewing) et internalisation}

\begin{definition}[Décalage des cotations (skewing)]\label{def:m2:fx-spot-microstructure:skew}
Le \emph{décalage des cotations (skewing)}\index{décalage des cotations (skewing)} consiste à déplacer le cours acheteur et le cours vendeur
d'un flux ensemble, par rapport à l'estimation du milieu par le fournisseur, pour attirer les transactions qui réduisent
son inventaire~: un fournisseur long en euros abaisse les deux prix pour que ses clients lui achètent
des euros et que moins d'entre eux lui en vendent (\cref{fig:m2:fx-spot-microstructure:skew}).
\end{definition}

Un gros fournisseur qui reçoit des milliers de demandes de clients par minute peut
internaliser une grande partie de son risque (livre~1, chapitre~10)~: une vente d'un
client est compensée par un achat d'un autre, et seul le solde doit être
couvert sur une plateforme. Le décalage des cotations accélère la compensation, au prix
d'une petite concession sur le prix, et chaque couverture que le fournisseur n'envoie pas
vers une plateforme est une information qu'il ne révèle pas. Les modèles d'inventaire qui sous-tendent
les skews, et qui arbitrent entre le spread espéré et le risque de détenir une position, sont ceux
du livre~1 et des jeux de teneurs de marché du
\cref{ch:m2:market-making-games}.

\begin{omfigure}
\centering
\begin{tikzpicture}[font=\small,y=1.1cm]
\draw[gray,dashed] (-0.3,0) node[left,font=\footnotesize] {mid du fournisseur} -- (9.6,0);
\foreach \x/\t/\s in {1.5/position à plat/0,5/{long de\\20~millions d'EUR}/-0.35,8.5/{short de\\20~millions d'EUR}/0.35} {
  \draw[omDef,very thick] (\x-0.6,{\s+0.5}) -- (\x+0.6,{\s+0.5}) node[right,font=\footnotesize] {ask};
  \draw[omThm,very thick] (\x-0.6,{\s-0.5}) -- (\x+0.6,{\s-0.5}) node[right,font=\footnotesize] {bid};
  \node[font=\footnotesize,align=center] at (\x,-1.45) {\t};
}
\draw[omProp,thick,-{Stealth}] (5,-0.05) -- (5,-0.32);
\draw[omProp,thick,-{Stealth}] (8.5,0.05) -- (8.5,0.32);
\end{tikzpicture}

\omcaption{Décalage des cotations. Sans position, le fournisseur cote de façon symétrique
autour de son milieu~; long en euros, il abaisse les deux prix, pour que les acheteurs viennent
à lui et que les vendeurs aillent ailleurs~; short, il les relève. La fourchette ne change pas,
c'est le centre qui se déplace. Schéma.}\label{fig:m2:fx-spot-microstructure:skew}
\end{omfigure}

\section{Le code de conduite mondial}

\begin{definition}[Code de change FX Global]\label{def:m2:fx-spot-microstructure:code}
Le \emph{FX Global Code}\index{FX Global Code} est un ensemble de principes de bonnes
pratiques pour le marché de gros du change, tenu à jour par le Global Foreign Exchange
Committee, qui réunit des banques centrales et des acteurs du secteur privé. Il n'a pas force de loi~:
les participants y adhèrent en signant une déclaration d'engagement.
\end{definition}

Le Code est apparu en 2017, deux ans après les affaires de sanction liées au
fixing du \cref{ch:m2:fixings-and-flows}, et le \omterm{def:m2:fx-spot-microstructure:lastlook}{last look} est l'une de ses sections les plus
discutées. Le principe~17, tel que mis à jour en
décembre 2024, exige que les fournisseurs qui utilisent le \omterm{def:m2:fx-spot-microstructure:lastlook}{last look} en soient transparents~:
qu'ils indiquent si et comment les variations de prix, dans un sens comme dans l'autre, influent sur la décision,
combien de temps il faut et pourquoi ils y recourent. Il définit le \omterm{def:m2:fx-spot-microstructure:lastlook}{last look} comme un contrôle de risque, destiné à
vérifier qu'une demande est valide et que son prix reste cohérent avec le
marché, et interdit de s'en servir pour collecter de l'information ou pour négocier sur la demande
pendant la fenêtre, que ce soit en ajustant les prix ou en couvrant la position. En août 2021,
le Comité a publié des modèles selon lesquels les fournisseurs divulguent leurs pratiques.

\begin{remark}[Ce que le Code laisse ouvert]\label{rem:m2:fx-spot-microstructure:open}
Le Code n'interdit pas les contrôles asymétriques et ne fixe pas de \omterm{def:m2:fx-spot-microstructure:lastlook}{temps de rétention}. Il demande
aux fournisseurs de les divulguer et aux clients de lire ces divulgations~; la pression
qui subsiste est concurrentielle. Les clients mesurent par fournisseur leurs taux d'exécution, leurs \omterm{def:m2:fx-spot-microstructure:lastlook}{temps de
rétention} et leurs \omterm{def:m2:fx-spot-microstructure:markout}{mark-outs}, comme le fait le tutoriel, et orientent leurs demandes
en conséquence.
\end{remark}

\begin{dated}{2026-09}{Le Code aujourd'hui}\label{dat:m2:fx-spot-microstructure:code}
Le texte en vigueur du \omterm{def:m2:fx-spot-microstructure:code}{FX Global Code} est daté de décembre 2024~; son principe~17 sur
le \omterm{def:m2:fx-spot-microstructure:lastlook}{last look} est tel que résumé ci-dessus. Le Global Foreign Exchange Committee, créé
en mai 2017 pour le tenir à jour, a publié en août 2021 un document d'orientation sur le
\omterm{def:m2:fx-spot-microstructure:lastlook}{last look} et des modèles standard selon lesquels les teneurs de liquidité divulguent leurs
pratiques. L'adhésion reste volontaire, par déclaration d'engagement.
\end{dated}

\section{Tutoriel~: analyse des coûts de transaction du last look}\label{tut:m2:fx-spot-microstructure:tca}

\begin{tutorial}
\textbf{Objectif.} Simuler un flux avec des clients informés et des clients non informés, appliquer
des contrôles de \omterm{def:m2:fx-spot-microstructure:lastlook}{last look} asymétriques et symétriques, et mesurer les \omterm{def:m2:fx-spot-microstructure:lastlook}{taux de rejet}, les taux d'exécution
et les \omterm{def:m2:fx-spot-microstructure:markout}{mark-outs} en fonction du \omterm{def:m2:fx-spot-microstructure:lastlook}{temps de rétention}. \textbf{État final~:}
\cref{fig:m2:fx-spot-microstructure:hold}, \cref{ex:m2:fx-spot-microstructure:window}
et les nombres du problème du week-end.
\begin{tutsteps}
\item \textbf{Le contrôle et le flux}~: le mouvement de chaque demande pendant
le \omterm{def:m2:fx-spot-microstructure:lastlook}{temps de rétention} et après.
\omcode{code/firm/lastlook/firm_lastlook.py}{24}{43}{Le contrôle de last look et un flux simulé de demandes.}\label{lst:m2:fx-spot-microstructure:sim}
\item \textbf{L'analyse}~: taux d'exécution, \omterm{def:m2:fx-spot-microstructure:lastlook}{taux de rejet} par type de client, \omterm{def:m2:fx-spot-microstructure:markout}{mark-out},
et la forme close de la \cref{prop:m2:fx-spot-microstructure:transfer}.
\omcode{code/firm/lastlook/firm_lastlook.py}{46}{69}{Analyse des coûts de transaction et transfert espéré.}\label{lst:m2:fx-spot-microstructure:tca}
\item \textbf{Exécuter} \texttt{lastlook\_demo.by\_hold} pour chaque politique,
\texttt{lastlook\_demo.window()} et \texttt{fig\_lastlook.py}.
\end{tutsteps}
\textbf{Pour aller plus loin.} Remplacer l'edge fixe des clients informés par un
edge aléatoire, et trouver le \omterm{def:m2:fx-spot-microstructure:lastlook}{temps de rétention} auquel un contrôle symétrique a rejeté
90~\% d'entre eux~; tracer ensuite le taux d'exécution des clients non informés à ce \omterm{def:m2:fx-spot-microstructure:lastlook}{temps de rétention}.
\end{tutorial}

\section{Construction~: l'analyseur de last look}\label{bld:m2:fx-spot-microstructure:lastlook}

\begin{build}
\bfield{Rôle} La petite société de trading prend et fournit de la liquidité au comptant sur le change~:
en preneur, elle doit savoir ce que coûte à chaque fournisseur son \omterm{def:m2:fx-spot-microstructure:lastlook}{last look}~; en fournisseur,
elle doit pouvoir montrer à ses clients que son propre contrôle est symétrique et bref.
\bfield{Interface} \texttt{check(move, threshold, policy)}~; \texttt{simulate(n, sigma, hold, threshold, policy, informed\_share, edge,
horizon, seed)} renvoyant des enregistrements \texttt{Request}~; \texttt{tca(requests,
half)}~; \texttt{expected\_transfer(sigma, hold, threshold, policy)}.
\bfield{Règles} Mouvements en points de base, positifs en faveur du client~;
milieu brownien~; politiques \texttt{none}, \texttt{asymmetric}, \texttt{symmetric}~;
simulation à graine fixée.
\bfield{Tests d'acceptation} \texttt{code/firm/lastlook/tests/}~: les définitions des
politiques~; aucun \omterm{def:m2:fx-spot-microstructure:lastlook}{last look} n'exécute tout~; le transfert asymétrique simulé à moins de
2~\% de la forme close~; les demandes informées rejetées beaucoup plus souvent que les
demandes non informées.
\bfield{Pour aller plus loin} Lire de vraies données d'exécution et de rejet (un journal FIX) et calculer
pour chaque fournisseur la distribution des \omterm{def:m2:fx-spot-microstructure:lastlook}{temps de rétention} et les \omterm{def:m2:fx-spot-microstructure:markout}{mark-outs}~; tester statistiquement si les
rejets d'un fournisseur dépendent du sens du mouvement.
\end{build}

\begin{omsources}
\item New York State Department of Financial Services, Consent Order, In the Matter of Barclays Bank PLC, novembre 2015.
\item Global Foreign Exchange Committee, \emph{FX Global Code}, mis à jour
décembre 2024~; communiqué de presse sur les orientations relatives au last look et aux modèles d'information,
août 2021.
\item Bank for International Settlements, \emph{Quarterly Review}, décembre
2022, encadré sur le marché au comptant fragmenté.
\end{omsources}

\section{Exercices}

\begin{exercise}[$\star$]\label{exo:m2:fx-spot-microstructure:1}
Un fournisseur vend à un client 1~million d'euros à 1.1464 lorsque le mid vaut 1.1463. Une
seconde plus tard, le mid vaut 1.1465. Donner le \omterm{def:m2:fx-spot-microstructure:markout}{mark-out} du fournisseur à une seconde en
dollars et en points de base.
\end{exercise}

\begin{exercise}[$\star$]\label{exo:m2:fx-spot-microstructure:2}
Donner l'écart type du mouvement de prix sur des \omterm{def:m2:fx-spot-microstructure:lastlook}{temps de rétention} de 25, 100 et 400
millisecondes, à raison de 0.01 point de base par racine carrée de milliseconde.
\end{exercise}

\begin{exercise}[$\star$]\label{exo:m2:fx-spot-microstructure:3}
Que doit divulguer un fournisseur qui recourt au \omterm{def:m2:fx-spot-microstructure:lastlook}{last look} au titre du principe~17, et que
ne doit-il pas faire pendant la fenêtre~?
\end{exercise}

\begin{exercise}[$\star\star$]\label{exo:m2:fx-spot-microstructure:4}
Donner le transfert par demande d'un contrôle asymétrique avec un \omterm{def:m2:fx-spot-microstructure:lastlook}{temps de rétention} de 25~millisecondes
et un seuil nul, puis le comparer à celui obtenu pour un \omterm{def:m2:fx-spot-microstructure:lastlook}{temps de rétention} de 100~millisecondes.
\end{exercise}

\begin{exercise}[$\star\star$]\label{exo:m2:fx-spot-microstructure:5}
Un fournisseur long de 20~millions d'euros décale sa \omterm{def:m2:fx-spot-microstructure:stream}{cotation en streaming} EURUSD de 0.1 pip. Dans quel sens,
et pourquoi cela peut-il être moins coûteux que de se couvrir sur une \omterm{def:m2:the-fx-market:venue}{plateforme principale}~?
\end{exercise}

\begin{exercise}[$\star\star$]\label{exo:m2:fx-spot-microstructure:6}
Pourquoi un contrôle symétrique rejette-t-il davantage de demandes non informées qu'un
contrôle asymétrique à seuil égal, tout en leur coûtant moins~?
\end{exercise}

\begin{exercise}[$\star\star\star$]\label{exo:m2:fx-spot-microstructure:7}
\emph{Programmation.} À partir de \texttt{by\_hold}, donner les \omterm{def:m2:fx-spot-microstructure:lastlook}{taux de rejet} des demandes informées et non informées
à 100~millisecondes selon chaque politique, ainsi que les \omterm{def:m2:fx-spot-microstructure:markout}{mark-outs}.
\end{exercise}

\begin{exercise}[$\star\star\star$]\label{exo:m2:fx-spot-microstructure:8}
\emph{Repérer l'erreur.} «~Notre \omterm{def:m2:fx-spot-microstructure:lastlook}{taux de rejet} n'est que de 5~\%, donc notre \omterm{def:m2:fx-spot-microstructure:lastlook}{last look} ne coûte
rien à nos clients.~» Corriger cette affirmation.
\end{exercise}

\section{Problème~: La fenêtre asymétrique}

\begin{problem}[{Problème du week-end --- ce qu'une fenêtre de last look prend à un client}]\label{pb:m2:fx-spot-microstructure:1}
Une société de gestion envoie chaque jour à un fournisseur 2~milliards de dollars de demandes EURUSD,
qui recourt au \omterm{def:m2:fx-spot-microstructure:lastlook}{last look} avec un \omterm{def:m2:fx-spot-microstructure:lastlook}{temps de rétention} de 100~millisecondes et un
seuil de 0.05 point de base. Le mid évolue de 0.01 point de base par racine carrée de milliseconde, le
demi-spread vaut 0.4 point de base, et le flux du gestionnaire est non informé.

\medskip\noindent\textbf{Partie I --- La fenêtre.}
\begin{enumerate}
\item Donner l'écart type du mouvement sur le \omterm{def:m2:fx-spot-microstructure:lastlook}{temps de rétention}.
\item Sous un contrôle asymétrique, quelle part des demandes est rejetée~?
\item Et sous un contrôle symétrique de même seuil~?
\item Donner le transfert par demande dans chacun des deux cas.
\item Donner le transfert asymétrique sans seuil.
\end{enumerate}

\medskip\noindent\textbf{Partie II --- En valeur.}
\begin{enumerate}[resume]
\item Donner le coût quotidien de la fenêtre asymétrique.
\item Donner ce coût sur 250 jours de négociation.
\item Comment se compare-t-il au demi-spread que paie le gestionnaire~?
\item Donner le coût avec un \omterm{def:m2:fx-spot-microstructure:lastlook}{temps de rétention} de 25~millisecondes et sans seuil.
\item Pourquoi le coût croît-il avec la racine carrée du \omterm{def:m2:fx-spot-microstructure:lastlook}{temps de rétention}~?
\end{enumerate}

\medskip\noindent\textbf{Partie III --- Ce que voit le gestionnaire.}
\begin{enumerate}[resume]
\item Quelles statistiques révèlent un contrôle asymétrique~?
\item Pourquoi un faible \omterm{def:m2:fx-spot-microstructure:lastlook}{taux de rejet} ne suffit-il pas~?
\item Que doit demander le gestionnaire au fournisseur, en citant le \omterm{def:m2:fx-spot-microstructure:code}{FX Global Code}~?
\item Que coûte encore un rejet, au-delà du transfert~?
\item Pourquoi le gestionnaire pourrait-il encore préférer un fournisseur recourant au \omterm{def:m2:fx-spot-microstructure:lastlook}{last look}~?
\end{enumerate}

\medskip\noindent\textbf{Partie IV --- Jugement.}
\begin{enumerate}[resume]
\item Que reprochait l'ordonnance de New York au contrôle de Barclays~?
\item Un contrôle symétrique est-il équitable pour le client~?
\item Comment une \omterm{def:m2:fx-spot-microstructure:stream}{cotation en streaming} à prix ferme, sans \omterm{def:m2:fx-spot-microstructure:lastlook}{last look}, se protège-t-elle à la place~?
\item Énoncer le \emph{résultat nommé}~: ce que la fenêtre asymétrique prend chaque année
au gestionnaire.
\item En une phrase~: que coûte une fenêtre de \omterm{def:m2:fx-spot-microstructure:lastlook}{last look} à un client~?
\end{enumerate}
\end{problem}

\section{Questions d'entretien}

\begin{interviewq}[$\star$ \iqroles{trader, développeur}]\label{iq:m2:fx-spot-microstructure:1}
Qu'est-ce que le \omterm{def:m2:fx-spot-microstructure:lastlook}{last look}, et pourquoi les fournisseurs de liquidité y ont-ils recours~?
\end{interviewq}

\begin{interviewq}[$\star$ \iqroles{chercheur, trader}]\label{iq:m2:fx-spot-microstructure:2}
Comment sauriez-vous, à partir de vos exécutions, si le \omterm{def:m2:fx-spot-microstructure:lastlook}{last look} d'un fournisseur est
asymétrique~?
\end{interviewq}

\begin{interviewq}[$\star\star$ \iqroles{chercheur}]\label{iq:m2:fx-spot-microstructure:3}
Qu'est-ce qu'un \omterm{def:m2:fx-spot-microstructure:markout}{mark-out}, et comment utilisez-vous les \omterm{def:m2:fx-spot-microstructure:markout}{mark-outs} pour répartir vos clients par paliers~?
\end{interviewq}

\begin{interviewq}[$\star\star$ \iqroles{trader}]\label{iq:m2:fx-spot-microstructure:4}
Vous êtes long de 50~millions d'euros issus du flux client. Décalez-vous vos cotations, couvrez-vous ou attendez-vous~?
\end{interviewq}

\begin{interviewq}[$\star\star$ \iqroles{banque, trader}]\label{iq:m2:fx-spot-microstructure:5}
Qu'exige le \omterm{def:m2:fx-spot-microstructure:code}{FX Global Code} en matière de \omterm{def:m2:fx-spot-microstructure:lastlook}{last look}, et que n'exige-t-il
pas~?
\end{interviewq}

\begin{interviewq}[$\star\star\star$ \iqroles{développeur}]\label{iq:m2:fx-spot-microstructure:6}
Concevez un moteur de prix qui diffuse la cotation EURUSD en streaming à un millier de clients répartis en plusieurs
paliers, qui applique un contrôle symétrique de \omterm{def:m2:fx-spot-microstructure:lastlook}{last look} dans un budget de latence de dizaines
de microsecondes.
\end{interviewq}
